% !TeX program = pdflatex
\documentclass[11pt,a4paper]{article}

% Compile with pdfLaTeX:
%   latexmk -pdf legged_mani_analysis_draft.tex
\usepackage[utf8]{inputenc}
\usepackage[T1]{fontenc}

\usepackage{amsmath,amssymb,amsthm}
\usepackage{bm}
\usepackage{geometry}
\usepackage{booktabs}
\usepackage{array}
\usepackage{algorithm}
\usepackage{algorithmic}
\usepackage{xcolor}
\usepackage{hyperref}

\geometry{margin=2.4cm}
\hypersetup{
  colorlinks=true,
  linkcolor=blue!55!black,
  urlcolor=blue!55!black,
  citecolor=blue!55!black
}
\setcounter{tocdepth}{2}

\newcommand{\R}{\mathbb{R}}
\newcommand{\norm}[1]{\left\lVert #1 \right\rVert}
\newcommand{\sqnorm}[1]{\left\lVert #1 \right\rVert^2}
\newcommand{\wnorm}[2]{\left\lVert #1 \right\rVert_{#2}^2}
\newcommand{\clip}{\operatorname{clip}}
\newcommand{\FK}{\operatorname{FK}}
\newcommand{\diag}{\operatorname{diag}}
\newcommand{\dist}{\operatorname{dist}}
\newcommand{\code}[1]{\texttt{#1}}

\title{\code{legged\_mani}: B2--Z1 Whole-Body MPPI\\
Mathematical Formulation and Implementation Analysis}
\author{Notes based on analysis of the \code{legged\_mani} codebase}
\date{\today}

\begin{document}
\maketitle
\tableofcontents
\newpage

%=============================================================================
\section{Scope and Executive Summary}
%=============================================================================

This document analyzes the \code{locomani} task in \code{legged\_mani}, a
whole-body Model Predictive Path Integral (MPPI) controller for a Unitree B2
quadruped equipped with a reduced four-degree-of-freedom Unitree Z1 arm. The
analysis is based primarily on the following files:
\begin{itemize}
  \item MPPI sampling and rollout:
        \code{control/controllers/base\_controller.py},
  \item locomani controller:
        \code{control/controllers/mppi\_locomani.py},
  \item task and solver parameters:
        \code{utils/tasks.py} and \code{configs/mppi\_locomani.yml},
  \item robot and rollout models:
        \code{models/b2\_z1\_4dof.xml} and \code{models/b2\_z1\_base.xml},
  \item closed-loop execution:
        \code{interface/simulator.py}.
\end{itemize}

The main architectural decisions are:
\begin{enumerate}
  \item The observation is the complete 45-dimensional MuJoCo position and
        velocity state.
  \item The optimized input is a vector of 16 desired joint positions, not
        joint torques or contact forces.
  \item Full multibody dynamics and contact are evaluated through parallel
        MuJoCo rollouts instead of an explicit reduced-order dynamics model.
  \item The leg nominal trajectory is supplied by a periodic gait file, while
        the arm nominal posture is generated by damped least-squares inverse
        kinematics (IK) for the current end-effector target.
  \item Each rollout is scored by body and joint tracking, virtual PD effort,
        and end-effector tracking; arm--torso clearance is enforced only by
        hard rollout rejection.
  \item A rollout that violates the hard arm--torso clearance threshold is
        removed completely from the MPPI weighted average.
\end{enumerate}

Consequently, this controller is structurally different from a centroidal
legged MPC that optimizes contact forces under a prescribed contact-mode
schedule. Contact forces are internal MuJoCo variables, and gait timing acts as
a kinematic prior rather than a hard contact constraint.

%=============================================================================
\section{Robot Model, State, and Input}
%=============================================================================

\subsection{Degrees of Freedom}

The MuJoCo model dimensions are
\begin{equation}
  n_q=23,\qquad n_v=22,\qquad n_u=16.
\end{equation}
The actuated coordinates consist of twelve B2 leg joints and the first four Z1
arm joints. Z1 joints 5 and 6 and the gripper are fixed in the reduced model.

The actuator and joint order is
\begin{align*}
\bm{q}_j = [&q_{\mathrm{FR,h}},q_{\mathrm{FR,t}},q_{\mathrm{FR,c}},
q_{\mathrm{FL,h}},q_{\mathrm{FL,t}},q_{\mathrm{FL,c}},\\
&q_{\mathrm{RR,h}},q_{\mathrm{RR,t}},q_{\mathrm{RR,c}},
q_{\mathrm{RL,h}},q_{\mathrm{RL,t}},q_{\mathrm{RL,c}},
q_{a1},q_{a2},q_{a3},q_{a4}]^\top ,
\end{align*}
where $\mathrm{h}$, $\mathrm{t}$, and $\mathrm{c}$ denote hip, thigh, and calf,
respectively.

\subsection{Observation State}

The controller receives the concatenation of MuJoCo \code{qpos} and
\code{qvel}:
\begin{equation}
  \boxed{
  \bm{x}=
  \begin{bmatrix}
    \bm{p}_b\\
    \bm{q}_b\\
    \bm{q}_j\\
    \bm{v}_b\\
    \bm{\omega}_b\\
    \dot{\bm{q}}_j
  \end{bmatrix}
  \in\R^{45}.}
\end{equation}
The components are:
\begin{itemize}
  \item $\bm{p}_b\in\R^3$: floating-base position in the world frame,
  \item $\bm{q}_b=[q_w,q_x,q_y,q_z]^\top\in\mathbb{S}^3$: base quaternion,
  \item $\bm{q}_j\in\R^{16}$: actuated leg and arm joint positions,
  \item $\bm{v}_b,\bm{\omega}_b\in\R^3$: base linear and angular velocities,
  \item $\dot{\bm{q}}_j\in\R^{16}$: actuated joint velocities.
\end{itemize}

The exact array layout is summarized below.
\begin{center}
\begin{tabular}{ccl}
\toprule
Indices & Dimension & Quantity\\
\midrule
$0{:}3$ & 3 & base position $\bm{p}_b$\\
$3{:}7$ & 4 & base quaternion $\bm{q}_b$\\
$7{:}23$ & 16 & actuated joint positions $\bm{q}_j$\\
$23{:}26$ & 3 & base linear velocity $\bm{v}_b$\\
$26{:}29$ & 3 & base angular velocity $\bm{\omega}_b$\\
$29{:}45$ & 16 & actuated joint velocities $\dot{\bm{q}}_j$\\
\bottomrule
\end{tabular}
\end{center}

The MuJoCo FULLPHYSICS rollout buffer contains
$[t,\bm{q}_{\mathrm{pos}}^\top,\bm{q}_{\mathrm{vel}}^\top]^\top\in\R^{46}$.
The leading time element is removed before cost evaluation, leaving the
45-dimensional state above.

\subsection{Control Input and Actuator Model}

The optimized control is the desired position of every actuated joint. Before
rollout, each sampled command is clipped componentwise to the actuator control
range:
\begin{equation}
  \boxed{\bm{u}=\bm{q}_{j,\mathrm{des}}\in\R^{16}},
  \qquad
  u_i^{\mathrm{clip}}
  =\operatorname{sat}_{[u_i^-,u_i^+]}(u_i).
\end{equation}
For a scalar $z$, the saturation operator used here is
\begin{equation}
  \operatorname{sat}_{[a,b]}(z)
  := \min\!\left(\max(z,a),b\right).
\end{equation}
After action clipping, each affine position actuator computes an unsaturated
PD force and then applies its actuator-force limits:
\begin{align}
  \tau_i^{\mathrm{unsat}}
  &=k_{p,i}(u_i^{\mathrm{clip}}-q_{j,i})-k_{d,i}\dot q_{j,i},\\
  \tau_i^{\mathrm{MJ}} &=
  \operatorname{sat}_{[\tau_i^-,\tau_i^+]}
  \left(\tau_i^{\mathrm{unsat}}\right).
\end{align}
The B2 hip actuators use $(k_p,k_d)=(100,5)$, the thigh and calf actuators
use $(300,8)$, and all four arm actuators use $(300,5)$. Joint-position
limits and actuator-force limits are read from the MJCF model and applied by
the sampler and MuJoCo, respectively. Thus action clipping constrains desired
joint positions, whereas torque saturation constrains the force generated by
the position actuator.

%=============================================================================
\section{Prediction Dynamics and Contact}
%=============================================================================

\subsection{Discrete MuJoCo Dynamics}

The controller does not implement a closed-form reduced model
$\dot{\bm{x}}=f(\bm{x},\bm{u})$. One prediction step is instead represented as
\begin{equation}
  \boxed{
  \bm{x}_{k+1}=F_{\mathrm{MJ}}
  (\bm{x}_k,\bm{u}_k;\bm{\theta}_{\mathrm{model}}),
  }
\end{equation}
where $F_{\mathrm{MJ}}$ is evaluated by \code{rollout.rollout}. The model
captures:
\begin{itemize}
  \item floating-base and articulated rigid-body dynamics,
  \item position-PD actuators and force saturation,
  \item gravity, joint damping, and armature,
  \item foot--ground collision, friction, and MuJoCo soft-contact dynamics.
\end{itemize}

The default integration interval is $\Delta t=0.01\,\mathrm{s}$ and the
horizon is $H=40$, giving
\begin{equation}
  T_H=H\Delta t=0.4\,\mathrm{s}.
\end{equation}

\subsection{Consistency Between Prediction and Simulation}

The MPPI controller and the closed-loop simulator load separate MuJoCo model
instances. During simulator initialization, the contact cone, override flag,
\code{o\_solref}, \code{o\_solimp}, friction, and margin are copied from the
rollout model to the simulator model. The locomani configuration uses a
pyramidal friction cone and
\begin{equation}
  \code{o\_solref}=[0.02,1.0].
\end{equation}
This synchronization reduces prediction--execution mismatch in contact
dynamics.
% 빼도 될듯?

\subsection{Implicit Contact Scheduling}

Contact force is not part of $\bm{u}$, and there are no explicit stance or
swing equality constraints. A periodic leg-joint trajectory encourages the
intended swing and stance motion, but actual contact is determined by MuJoCo
collision detection. The gait scheduler is therefore a phase-aligned
kinematic prior, not a strict contact-mode scheduler.

%=============================================================================
\section{Finite-Horizon Stochastic Optimal Control Problem}
%=============================================================================

At each control update, the intended finite-horizon problem can be written as
\begin{equation}
\boxed{
\begin{aligned}
\min_{\bm{U}}\quad
  &J(\bm{U};\bm{x}_0)
  =\sum_{k=0}^{H-1}\ell_k(\bm{x}_{k+1},\bm{u}_k)
   +\Phi_{\mathrm{EE}}(\bm{x}_H),\\
\mathrm{s.t.}\quad
  &\bm{x}_{k+1}=F_{\mathrm{MJ}}(\bm{x}_k,\bm{u}_k),\\
  &\bm{u}_{\min}\leq\bm{u}_k\leq\bm{u}_{\max},\\
  &d_{k,s}\geq d_{\mathrm{hard}},\qquad s=1,2,3,\\
  &\bm{x}_0=\hat{\bm{x}}(t),
\end{aligned}}
\label{eq:finite_ocp}
\end{equation}
where $\bm{U}=[\bm{u}_0,\ldots,\bm{u}_{H-1}]$ and $d_{k,s}$ is the signed
clearance between arm capsule $s$ and the torso box.

The nonlinear, contact-rich problem in Eq.~\eqref{eq:finite_ocp} is not solved
with gradients. Instead, $N=30$ control sequences are sampled around the
current nominal trajectory, evaluated through parallel MuJoCo rollouts, and
combined by exponential weighting.

%=============================================================================
\section{Nominal Trajectory Construction}
%=============================================================================

\subsection{Leg Gait Reference}

The default \code{in\_place} gait file has the form
\begin{equation}
  \bm{G}\in\R^{32\times L},\qquad L=100,
\end{equation}
with rows ordered as
\begin{equation}
  \bm{G}=
  \begin{bmatrix}
    \bm{q}_{\ell}^{g}\\
    \bm{q}_{a}^{g}\\
    \dot{\bm{q}}_{\ell}^{g}\\
    \dot{\bm{q}}_{a}^{g}
  \end{bmatrix},
  \qquad
  \bm{q}_{\ell}^{g},\dot{\bm{q}}_{\ell}^{g}\in\R^{12\times L}.
\end{equation}
At every controller update, the gait scheduler circularly shifts the column
indices by one sample and exposes the next $H$ reference columns.

\subsection{Arm IK Reference}

For the active world-frame EE target $\bm{p}_e^*$, the controller solves a
position-only damped least-squares IK problem. At iteration $r$,
\begin{align}
  \bm{e}^{(r)}
  &=\bm{p}_e^*-\bm{p}_e(\bm{q}_a^{(r)}),\\
  \Delta\bm{q}_a^{(r)}
  &=\bm{J}_a^\top
  \left(\bm{J}_a\bm{J}_a^\top+\lambda_{\mathrm{IK}}^2\bm{I}\right)^{-1}
  \bm{e}^{(r)},\\
  \bm{q}_a^{(r+1)}
  &=\clip\left(
  \bm{q}_a^{(r)}+\alpha_{\mathrm{IK}}\Delta\bm{q}_a^{(r)}
  \right).
\end{align}
The default damping and step size are
$\lambda_{\mathrm{IK}}=0.05$ and $\alpha_{\mathrm{IK}}=0.7$, with at most 100
iterations.

IK is recomputed from the current posture at every MPPI update. To avoid
abrupt changes in the control prior, the new solution is filtered with a
smoothing factor of $0.3$ and a per-update joint change limit of
$0.04\,\mathrm{rad}$. IK only supplies a nominal arm posture; actual EE error
is evaluated directly along every rollout.

\subsection{Whole-Body Joint Reference and Residual Warm Start}

The controller replaces the arm rows of the gait file with the current IK
posture:
\begin{equation}
  \bm{r}_{j,k}=
  \begin{bmatrix}
    \bm{q}^{g}_{\ell,k}\\
    \bm{q}^{\mathrm{IK}}_a\\
    \dot{\bm{q}}^{g}_{\ell,k}\\
    \bm{0}_4
  \end{bmatrix}
  \in\R^{32}.
\end{equation}
Its position component defines the nominal action:
\begin{equation}
  \bar{\bm{u}}_k=
  \begin{bmatrix}
    \bm{q}^{g}_{\ell,k}\\
    \bm{q}^{\mathrm{IK}}_a
  \end{bmatrix}
  +\Delta\bm{u}^{\mathrm{warm}}_k.
\end{equation}

After optimization, the difference between the selected action and the
gait/IK nominal is stored as a residual and shifted by one step. The new
horizon tail is initialized with zero residual, so the newly exposed future
sample begins at the original gait/IK nominal. During the first 50 controller
updates, the nominal is linearly blended from the initial standing control to
the gait/IK reference.

%=============================================================================
\section{Sampling and MPPI Update}
%=============================================================================

\subsection{Rollout Sampling Modes}

The current controller exposes three main rollout modes:
\code{original}, \code{original\_spline}, and \code{safe\_spline}. Both the
controller constructor and \code{simulate\_mppi.py} select
\code{original\_spline} by default. All three modes use $K=4$ horizon-aligned
knots
\begin{equation}
  \mathcal{K}=\{\kappa_j\}_{j=0}^{K-1},
  \qquad
  \kappa_j=
  \operatorname{round}\left(\frac{j(H-1)}{K-1}\right),
\end{equation}
For $H=40$, this gives $\mathcal{K}=\{0,13,26,39\}$. The implementation
requires the rounded indices to be unique. Let
$\mathcal{S}_{\mathcal{K}}[\cdot]_k$ denote cubic-spline reconstruction at
horizon index $k$ from values defined at these knots. Each sample draws
\begin{equation}
  \bm{\epsilon}_{n,j}
  \sim\mathcal{N}(\bm{0},\diag(\bm{\sigma}^2)),
  \qquad j=0,\ldots,K-1.
\end{equation}

For the gait-aware modes, define the full-rate gait/IK nominal
$\bm{g}_k$, the shifted MPPI correction $\bm{c}_k$, and the clipped absolute
warm start
\begin{equation}
  \bm{a}_k
  =\clip(\bm{g}_k+\bm{c}_k,\bm{u}_{\min},\bm{u}_{\max}).
\end{equation}
The three sampling rules are as follows.

\paragraph{\code{original}.}
This mode disables the explicit gait/IK nominal
(\code{nominal\_from\_gait}$=$\code{False}) and maps internally to
\path{sample_type=cubic_original}. Here $\bm{a}_k$ denotes the
previous optimized absolute action sequence after receding-horizon shifting.
The knot actions are perturbed and the complete absolute action is
reconstructed:
\begin{equation}
  \boxed{
  \bm{u}_{n,k}^{\mathrm{original}}
  =\clip\left(
    \mathcal{S}_{\mathcal{K}}
    [\bm{a}_{\kappa_j}+\bm{\epsilon}_{n,j}]_k,
    \bm{u}_{\min},\bm{u}_{\max}
  \right).
  }
\end{equation}
After the MPPI update, the selected absolute sequence is shifted and its last
action is repeated at the new horizon tail.

\paragraph{\code{original\_spline}.}
This is the default mode and maps internally to
\path{sample_type=cubic_gait_absolute}. It enables the gait/IK nominal and
forms $\bm{a}_k$ from the full gait/IK/residual warm start, but then splines
that complete absolute trajectory:
\begin{equation}
  \boxed{
  \bm{u}_{n,k}^{\mathrm{original\_spline}}
  =\clip\left(
    \mathcal{S}_{\mathcal{K}}
    [\bm{a}_{\kappa_j}+\bm{\epsilon}_{n,j}]_k,
    \bm{u}_{\min},\bm{u}_{\max}
  \right).
  }
\end{equation}
Consequently, both the gait/IK component and the warm-started correction are
reconstructed from only four knots. Short full-rate gait features between
knots can therefore be smoothed.

\paragraph{\code{safe\_spline}.}
This mode maps internally to \path{sample_type=cubic_gait_residual}. It
preserves the complete full-rate gait/IK nominal and applies the cubic spline
only to the correction plus exploration noise:
\begin{equation}
  \boxed{
  \bm{u}_{n,k}^{\mathrm{safe\_spline}}
  =\clip\left(
    \bm{g}_k+
    \mathcal{S}_{\mathcal{K}}
    [\bm{c}_{\kappa_j}+\bm{\epsilon}_{n,j}]_k,
    \bm{u}_{\min},\bm{u}_{\max}
  \right).
  }
\end{equation}
It is therefore the only one of the three modes that preserves the nominal
gait/IK samples at every horizon step before final action clipping.

Candidate zero is overwritten with a mode-consistent noise-free baseline:
\begin{center}
\begin{tabular}{p{0.22\textwidth}p{0.62\textwidth}}
\toprule
Mode & Candidate-zero baseline\\
\midrule
\code{original}
& $\bm{u}_{0,k}=\clip(\bm{a}_k)$, the full shifted absolute warm start\\
\code{original\_spline}
& $\bm{u}_{0,k}=\clip(
   \mathcal{S}_{\mathcal{K}}[\bm{a}_{\kappa_j}]_k)$\\
\code{safe\_spline}
& $\bm{u}_{0,k}=\clip(
   \bm{g}_k+\mathcal{S}_{\mathcal{K}}[\bm{c}_{\kappa_j}]_k)$\\
\bottomrule
\end{tabular}
\end{center}
Thus every batch contains a zero-noise reference appropriate to the selected
sampling family. In \code{original\_spline} and \code{safe\_spline}, the
selected correction relative to $\bm{g}_k$ is shifted after optimization and
the new tail correction is set to zero.

The exploration standard deviations are $0.03$ for each hip, $0.10$ for each
thigh and calf, and $0.06\,\mathrm{rad}$ for each arm joint. The default
\code{in\_place} gait applies noise scale $1.0$.

\subsection{Cost Normalization and Weights}

Let $\mathcal{V}$ be the set of candidates that pass the hard collision test,
and let $S_n$ be the total rollout cost. The implementation normalizes costs by
the valid-sample range:
\begin{align}
  S_{\min}&=\min_{n\in\mathcal{V}}S_n,\\
  \Delta S&=\max_{n\in\mathcal{V}}S_n-S_{\min},\\
  w_n&=
  \begin{cases}
    \exp\!\left[
      -\dfrac{1}{\lambda_{\mathrm{MPPI}}}
      \dfrac{S_n-S_{\min}}{\Delta S}
    \right],
    &n\in\mathcal{V},\ \Delta S>0,\\
    1,&n\in\mathcal{V},\ \Delta S\simeq0,\\
    0,&n\notin\mathcal{V}.
  \end{cases}
\end{align}
The default temperature is $\lambda_{\mathrm{MPPI}}=0.1$.

The new sequence is the weighted average of the absolute candidate actions:
\begin{equation}
  \boxed{
  \bm{u}^{+}_k=
  \clip\left(
    \frac{\sum_{n=1}^{N}w_n\bm{u}_{n,k}}
         {\sum_{n=1}^{N}w_n+10^{-10}}
  \right).
  }
\end{equation}
Only $\bm{u}^{+}_0$ is executed; the remainder forms the next warm start.

If every candidate violates the hard clearance threshold, no weighted average
is formed. The last accepted safe trajectory is shifted by one step, and its
final action is repeated at the tail.

%=============================================================================
\section{Cost Function}
%=============================================================================

\subsection{State Reference}

The horizon state reference is assembled as
\begin{equation}
  \bm{x}^{\mathrm{ref}}_k=
  \begin{bmatrix}
    \bm{p}^{\mathrm{ref}}_b\\
    \bm{q}^{\mathrm{ref}}_b\\
    \bm{q}^{\mathrm{ref}}_{j,k}\\
    \bm{v}^{\mathrm{ref}}_b\\
    \bm{\omega}^{\mathrm{ref}}_b\\
    \dot{\bm{q}}^{\mathrm{ref}}_{j,k}
  \end{bmatrix}.
\end{equation}
For locomani, the body pose reference is the \code{stand} keyframe and the
base velocity reference is zero. Before cost evaluation, the current
world-frame base linear velocity is transformed into the body-local frame.
The joint references come from the gait/IK sequence defined above.

\subsection{Body and Joint State Cost}

The controller does not penalize a direct four-component quaternion
difference. It converts the current and reference base quaternions to roll and
pitch and overwrites the orientation slots of the error vector. Conceptually,
\begin{equation}
  \bm{e}_{x,k}=
  \begin{bmatrix}
    \bm{p}_{b,k}-\bm{p}^{\mathrm{ref}}_b\\
    e_{\mathrm{roll},k}\\
    e_{\mathrm{pitch},k}\\
    0\\
    0\\
    \bm{q}_{j,k}-\bm{q}^{\mathrm{ref}}_{j,k}\\
    \bm{v}_{b,k}-\bm{v}^{\mathrm{ref}}_b\\
    \bm{\omega}_{b,k}-\bm{\omega}^{\mathrm{ref}}_b\\
    \dot{\bm{q}}_{j,k}-\dot{\bm{q}}^{\mathrm{ref}}_{j,k}
  \end{bmatrix},
\end{equation}
and
\begin{equation}
  \ell_{x,k}=\bm{e}_{x,k}^{\top}\bm{Q}\bm{e}_{x,k}.
\end{equation}

The main diagonal state weights are:
\begin{itemize}
  \item base height: $1000$,
  \item base roll and pitch: $5000$ each,
  \item all 16 joint positions: $1200$ each,
  \item each base linear and angular velocity component: $0.1$.
\end{itemize}
The current weights for base $x/y$, yaw, and joint-velocity state errors are
zero.

\subsection{Virtual PD Effort Cost}

The control penalty is applied to a virtual PD effort rather than directly to
the desired joint positions:
\begin{equation}
  \tilde{\bm{\tau}}_k=
  \bm{K}_p(\bm{u}_k-\bm{q}_{j,k})
  -\bm{K}_d
  \dot{\bm{q}}_{j,k},
\end{equation}
\begin{equation}
  \ell_{u,k}=
  \tilde{\bm{\tau}}_k^\top\bm{R}\tilde{\bm{\tau}}_k,
  \qquad
  \bm{R}=10^{-3}\bm{I}_{16}.
\end{equation}
This is a cost-side virtual effort. It is not exactly the actuator force
applied by MuJoCo because the expression does not include force saturation.

\subsection{End-Effector Pose Cost}

Let $\bm{p}_{e,k}$ and $\bm{q}_{e,k}$ be the EE position and quaternion
obtained from MuJoCo FK or rollout sensors. The position cost is
\begin{equation}
  \ell_{p,k}=
  w_p\sqnorm{\bm{p}_{e,k}-\bm{p}_e^*},
  \qquad w_p=5000.
\end{equation}
The quaternion double cover is handled using the absolute quaternion dot
product:
\begin{align}
  e_{R,k}
  &=2\arccos\left[
    \clip\left(
      \left|\bm{q}_{e,k}^\top\bm{q}_e^*\right|,-1,1
    \right)
  \right],\\
  \ell_{R,k}
  &=w_R e_{R,k}^2.
\end{align}
The current configuration sets $w_R=0$, so EE orientation tracking is
disabled.

An additional terminal multiplier $\gamma_T=10$ is applied to the final EE
cost:
\begin{equation}
  \Phi_{\mathrm{EE}}(\bm{x}_H)
  =\gamma_T(\ell_{p,H-1}+\ell_{R,H-1}).
\end{equation}
Because the same error is already present in the final stage cost, its total
EE coefficient at the final step is $1+\gamma_T=11$.

\subsection{Arm--Torso Collision Model}

\subsubsection{Geometry}

Four frames---shoulder, elbow, wrist, and EE---define three arm capsules. Their
radii are
\begin{equation}
  \bm{r}_{\mathrm{cap}}=
  [0.030,0.030,0.0375]\;\mathrm{m}.
\end{equation}
The torso is approximated by the oriented box associated with the
\code{base\_collision} geometry, whose half-size is
\begin{equation}
  \bm{h}_b=[0.25,0.14,0.075]\;\mathrm{m}.
\end{equation}
The first $0.07\,\mathrm{m}$ of the proximal capsule is excluded to prevent
the structural shoulder mount from producing a false collision.

For capsule $s$, the signed clearance $d_{k,s}$ is the distance between the
capsule centerline and the torso box minus the capsule radius. Thus $d>0$
means separation, $d=0$ means contact, and $d<0$ means penetration.

\subsubsection{Hard Rejection Only}

The EE-tracking objective contains no soft arm--torso clearance cost. A
candidate is marked invalid if any state in its horizon violates
\begin{equation}
  d_{k,s}\geq d_{\mathrm{hard}}=0.005\,\mathrm{m}.
\end{equation}
Invalid candidates receive zero MPPI weight and therefore cannot contribute
to the updated action trajectory. The configured safe distance is retained
only for compatibility with the shared collision helper; it does not enter
the EE-tracking objective.

For efficiency, each arm segment is first sampled at nine points. Clearance
bounds classify segments that are definitely above or below the hard
threshold, and only ambiguous cases receive an exact
segment--axis-aligned-box distance calculation in the box frame.

\subsection{Total Stage and Rollout Cost}

The complete stage cost contains no collision penalty:
\begin{equation}
  \ell_k=
  \ell_{x,k}+\ell_{u,k}+\ell_{p,k}+\ell_{R,k}.
\end{equation}
The total cost of candidate rollout $n$ is
\begin{equation}
  \boxed{
  S_n=
  \sum_{k=0}^{H-1}\ell_k^{(n)}
  +\gamma_T\left(
    \ell_{p,H-1}^{(n)}+\ell_{R,H-1}^{(n)}
  \right).
  }
\end{equation}

%=============================================================================
\section{End-Effector Waypoints and Task Progression}
%=============================================================================

The locomani task tracks a sequence of world-frame EE positions:
\begin{equation}
\begin{split}
  \bm{p}_{e,0}^{*}&=[0.85,0.00,0.80]^\top,\\
  \bm{p}_{e,1}^{*}&=[2.00,0.15,0.80]^\top,\\
  \bm{p}_{e,2}^{*}&=[0.85,0.15,0.48]^\top,\\
  \bm{p}_{e,3}^{*}&=[0.00,0.00,0.48]^\top,\\
  \bm{p}_{e,4}^{*}&=[-1.00,0.00,0.80]^\top.
\end{split}
\end{equation}
Some targets are outside the fixed-base reach of the arm. Their world-frame
EE cost can therefore induce base locomotion in addition to arm motion.

The goal check is
\begin{equation}
  \norm{\bm{p}_e-\bm{p}_e^*}\leq0.03\,\mathrm{m},
  \qquad
  1-\left|\bm{q}_e^\top\bm{q}_e^*\right|\leq1.0.
\end{equation}
Since the absolute quaternion dot product lies in $[0,1]$, the configured
orientation threshold is always satisfied. In practice, waypoint completion
depends only on EE position. The waypoint timer advances only when the
simulator calls \code{next\_goal()} after a successful goal check.

% 최종 OCP 구성 추가

%=============================================================================
\section{Complete Receding-Horizon Algorithm}
%=============================================================================

\begin{algorithm}[ht]
\caption{\code{legged\_mani} locomani MPPI update}
\label{alg:locomani_mppi}
\begin{algorithmic}[1]
  \REQUIRE Current observation $\hat{\bm{x}}$, gait phase, previous safe plan
  \STATE Update arm IK for the active EE waypoint.
  \STATE For gait-aware modes, combine the leg gait, arm IK, and shifted
         correction to refresh the warm start.
  \STATE Generate $N=30$ candidates using the selected
         \code{original}, \code{original\_spline}, or
         \code{safe\_spline} rule.
  \STATE Replace candidate zero with the corresponding mode-consistent
         noise-free baseline.
  \STATE Roll out all candidates for $H=40$ steps in MuJoCo.
  \STATE Evaluate state, effort, and EE costs $S_n$, together with hard
         collision validity.
  \STATE Assign zero weight to candidates that violate hard clearance.
  \IF{at least one valid candidate exists}
    \STATE Normalize valid costs and compute exponential weights.
    \STATE Compute $\bm{U}^{+}$ as the weighted mean of absolute candidate
           actions.
    \STATE Store $\bm{U}^{+}$ as the most recent safe plan.
  \ELSE
    \STATE Shift and reuse the previous safe plan as $\bm{U}^{+}$.
  \ENDIF
  \STATE Shift either the absolute plan (\code{original}) or the optimized
         gait/IK correction (gait-aware modes), and advance gait phase.
  \RETURN First desired joint-position command $\bm{u}^{+}_0$
\end{algorithmic}
\end{algorithm}

The default closed-loop simulator invokes this update at $100\,\mathrm{Hz}$
and applies the first action for one MuJoCo step. It then evaluates the goal
condition and updates the waypoint manager if the active target was reached.

%=============================================================================
\section{Default Configuration Summary}
%=============================================================================

\begin{center}
\begin{tabular}{>{\raggedright\arraybackslash}p{0.36\textwidth}
                >{\raggedright\arraybackslash}p{0.47\textwidth}}
\toprule
Parameter & Default value\\
\midrule
Robot & Unitree B2 with reduced four-DoF Z1 arm\\
State / input dimension & $45$ / $16$\\
Input interpretation & desired joint position\\
Time step / horizon & $0.01\,\mathrm{s}$ / 40 steps ($0.4\,\mathrm{s}$)\\
Samples / workers & 30 / 5\\
Temperature & $0.1$\\
Rollout mode / knots & \code{original\_spline} (default) / four\\
Initial keyframe & \code{stand}\\
Default gait & \code{in\_place}, 100 samples per cycle\\
Rollout contact & pyramidal cone, \code{o\_solref}$=[0.02,1.0]$\\
EE position/orientation weights & 5000 / 0\\
EE terminal scale & 10\\
Collision hard distance & $0.005\,\mathrm{m}$\\
Collision soft weight & 0 (not included in the objective)\\
\bottomrule
\end{tabular}
\end{center}

%=============================================================================
\section{Comparison with OCS2 Legged MPC}
%=============================================================================

\begin{center}
\begin{tabular}{p{0.22\textwidth}p{0.33\textwidth}p{0.35\textwidth}}
\toprule
Aspect & OCS2 legged MPC & \code{legged\_mani} locomani\\
\midrule
Dynamics & Centroidal/SRBD or full centroidal & MuJoCo full multibody rollout\\
Optimized input & Contact forces and joint velocities & Desired joint positions\\
Contact transitions & Prescribed mode schedule & Gait prior and MuJoCo collision\\
Constraint handling & Equalities, inequalities, barriers & Action clipping and rollout rejection\\
Optimization method & SQP, DDP, or IPM & Sampling-based MPPI-style update\\
Derivative information & Dynamics and cost derivatives & Derivative-free rollout scoring\\
Manipulation objective & Not part of the basic legged example & EE world-pose cost and arm IK prior\\
Collision avoidance & Added as problem-specific constraints & Three arm capsules against torso OBB\\
\bottomrule
\end{tabular}
\end{center}

%=============================================================================
\section{Implementation Caveats and Limitations}
%=============================================================================

\begin{enumerate}
  \item \textbf{Difference from canonical MPPI.}
  The implementation normalizes cost by the batch min--max range and averages
  absolute candidate actions instead of applying a weighted noise increment
  with a control-noise correction term. Consequently, $\lambda=0.1$ controls
  selection pressure over normalized relative cost, not temperature in the
  original physical cost units.

  \item \textbf{EE orientation is disabled.}
  The quaternion cost is zero and the orientation success threshold is
  nonrestrictive. The current task is an EE position task. Enforcing a general
  six-dimensional pose with a four-DoF arm would require a reachable-task
  analysis and different objective design.

  \item \textbf{The collision model is incomplete.}
  Hard rejection covers only three arm capsules against one torso box. It does
  not guarantee arm--leg, arm--ground, gripper--environment, or general
  self-collision safety.

  \item \textbf{Gait timing is not a contact constraint.}
  If gait tracking becomes weak, the desired stance/swing timing is not
  guaranteed. If joint tracking is too strong, it can suppress the whole-body
  motion needed to reduce EE error.

  \item \textbf{Cost-side effort differs from applied actuator force.}
  The virtual PD effort omits actuator-force saturation. Large tracking errors
  can therefore produce a discrepancy between the effort penalty and the
  force actually applied in MuJoCo.

  \item \textbf{Real-time performance must be measured.}
  Every update evaluates $30\times40$ contact-dynamics steps and rollout
  sensors. Sustained $100\,\mathrm{Hz}$ operation depends on the target CPU and
  should be verified from wall-clock timing rather than the nominal simulation
  time step.

  \item \textbf{Waypoint dwell is not necessarily continuous.}
  The waypoint timer increments only on successful goal checks and is not
  reset when the EE leaves the threshold. It therefore counts successful
  invocations rather than enforcing a strict continuous dwell interval.
\end{enumerate}

%=============================================================================
\section{Recommended Validation}
%=============================================================================

The following measurements are useful for connecting the formulation to
closed-loop performance:
\begin{enumerate}
  \item mean, 95th, and 99th percentile update time and the 10 ms deadline-miss
        rate,
  \item fraction of valid rollouts, minimum torso clearance, and the number of
        all-invalid fallbacks,
  \item EE position error and waypoint transition time,
  \item base height, roll/pitch, foot contacts, and slip,
  \item gait/IK nominal, selected actions, and warm-start residual magnitude,
  \item actuator saturation rate and virtual PD effort,
  \item one-step error between the predicted rollout and the executed
        simulator state.
\end{enumerate}

Useful ablations include: hard collision rejection versus collision disabled,
hard-clearance threshold sensitivity, zero arm exploration noise, removal of
the terminal EE multiplier, and removal of gait-residual warm starting.

%=============================================================================
\section{Pipeline Summary}
%=============================================================================

\begin{equation*}
\boxed{
\begin{array}{c}
\text{Current MuJoCo state }\hat{\bm{x}}_t\\
\downarrow\\
\text{Phase-aligned leg gait + arm IK for the active EE target}\\
\downarrow\\
\text{Whole-body nominal action + warm-started residual}\\
\downarrow\\
\text{Mode-selected cubic-spline sampling }(30\times40\times16)\\
\downarrow\\
\text{Parallel MuJoCo rollout + EE/FK sensors}\\
\downarrow\\
\text{State/effort + EE cost and hard-collision validity}\\
\downarrow\\
\text{Hard-collision rejection + exponential weighting}\\
\downarrow\\
\text{Execute the first desired joint-position action}\\
\downarrow\\
\text{Shift gait phase and residual for the next MPC update}
\end{array}}
\end{equation*}

The defining design choice is to sample desired joint-position trajectories
through full MuJoCo dynamics rather than optimizing contact forces over a
reduced-order model. The gait and IK modules provide a structured prior, while
the MPPI update searches for whole-body residuals that balance stability, EE
tracking, virtual effort, and torso clearance.

\end{document}
